\documentclass[11pt]{article}

\usepackage{amsmath}
\usepackage{amssymb}
\usepackage{hyperref}
\usepackage[margin=1in]{geometry}
\usepackage{tabularx}
\usepackage{ragged2e}
\usepackage{url}

\title{Notes: Off-policy strategy}
\date{}

\begin{document}
\maketitle

The on-policy algorithm (\texttt{on\_policy\_mc\_control}, see
\href{../../on_policy_monte_carlo/opmc_c/Section_4_on_policy_constant_alpha_mc.ipynb}{\path{rl_course/on_policy_monte_carlo/opmc_c/Section_4_on_policy_constant_alpha_mc.ipynb}})
uses the same $\epsilon$-greedy policy both to \emph{generate} the episodes (the \textbf{behavior policy}) and to be \emph{improved} (the \textbf{target policy}). This keeps things simple but forces a compromise --- the policy must keep exploring ($\epsilon > 0$) even while we'd ultimately like it to be greedy/deterministic.

\textbf{Off-policy} methods break this coupling by using two distinct policies:

\begin{itemize}
    \item \textbf{Behavior policy} $b(a|s)$: generates the experience (the actions actually taken in the environment). It must keep exploring, so it's typically $\epsilon$-soft (e.g. $\epsilon$-greedy).
    \item \textbf{Target policy} $\pi(a|s)$: the policy we are actually trying to learn/evaluate. It can be fully greedy with respect to $Q(s,a)$, since it never has to be executed during data collection.
\end{itemize}

This requires $b$ to have \textbf{coverage} of $\pi$: whenever $\pi(a|s) > 0$, we need $b(a|s) > 0$ as well, otherwise we could never observe the actions $\pi$ would take.

\section*{Why we need importance sampling}

Returns $G_t$ collected under $b$ are distributed according to $b$, not $\pi$. To use them to estimate $Q_\pi(s,a)$, each return has to be reweighted by the \textbf{importance-sampling ratio}, i.e. the relative probability of the trajectory under $\pi$ vs. under $b$:

\begin{equation}
\rho_{t:T-1} = \prod_{k=t}^{T-1} \frac{\pi(A_k \mid S_k)}{b(A_k \mid S_k)}
\end{equation}

Two common estimators built from this ratio:

\begin{itemize}
    \item \textbf{Ordinary importance sampling}: $Q(s,a) = \frac{1}{N}\sum_i \rho_i G_i$ --- unbiased, but can have very high (even unbounded) variance since $\rho$ is a product of many ratios.
    \item \textbf{Weighted importance sampling}: $Q(s,a) = \frac{\sum_i \rho_i G_i}{\sum_i \rho_i}$ --- biased (bias vanishes asymptotically), but much lower variance in practice, so it's generally preferred.
\end{itemize}

\section*{A common special case: greedy target policy}

If $\pi$ is chosen to be the \textbf{greedy} policy w.r.t. $Q$ (deterministic), then $\pi(A_k|S_k)$ is either $1$ (the greedy action was taken) or $0$ (a different, exploratory action was taken). This means:

\begin{itemize}
    \item Any time the behavior policy takes a non-greedy (exploratory) action, $\rho$ becomes $0$ for the rest of that episode's prefix, and the trailing part of the trajectory can be discarded for the update.
    \item This is exactly the structure exploited by algorithms like \textbf{Q-learning}, which can be seen as an (incremental, bootstrapped) off-policy method with an implicit greedy target policy --- no explicit importance-sampling ratio is needed there because the TD target already uses $\max_a Q(s', a)$.
\end{itemize}

\section*{On-policy vs. off-policy at a glance}

\begin{center}
\begin{tabularx}{\textwidth}{|l|X|X|}
\hline
 & \textbf{On-policy} & \textbf{Off-policy} \\
\hline
Behavior policy & Same as target & Different from target (must have coverage) \\
\hline
Target policy & Must stay exploratory ($\epsilon$-soft) & Can be fully greedy/deterministic \\
\hline
Sample reweighting & Not needed & Needed (importance sampling) \\
\hline
Variance & Lower & Higher (especially ordinary IS) \\
\hline
Data reuse & Harder --- policy keeps changing & Easier --- can learn from old/other agents' data, human demonstrations, replay buffers \\
\hline
\end{tabularx}
\end{center}

Off-policy learning is what makes it possible to learn a target policy from data generated by a \emph{different} source entirely (an older version of the policy, a hand-coded exploratory policy, logged human data, etc.), which is central to methods like \textbf{Q-learning}, \textbf{Expected SARSA} (off-policy variant), and off-policy \textbf{actor-critic} / \textbf{deep RL} algorithms (e.g. DQN, using a replay buffer of past experience).

\section*{Resources}

\begin{enumerate}
    \item[{[1]}] \href{https://web.stanford.edu/class/psych209/Readings/SuttonBartoIPRLBook2ndEd.pdf}{Reinforcement Learning: An Introduction. Ch. 5: Monte Carlo Methods}
\end{enumerate}

\end{document}
